\scanpage{21}
\mspage{\handwritten{12}}
$H^i(Y^{\deleted{\unsure{it}}}_{\handwritten{\bar t}},\underline{Q}_{\handwritten{\unsure{\ell|l}}})$, qui peut s'écrire, grâce au\deleted{\unsure{é|x}} théorè\unsure{\deleted{p}m|\deleted{n}m|m}e\deleted{s} de Lefschetz
\margin{\handwritten{faible \unsure{e}}}
\[
\begin{aligned}
&\deleted{2}\handwritten{\sum_{0\leq i\leq n-2}}\inserted{\handwritten{(-1)^i}}\operatorname{Trace}\operatorname{Frob}_{H^i(X_{\bar t},Q_{\handwritten{\ell}})}\deleted{(-1)^i}\inserted{\handwritten{(1+q^{2n-1-i})}}\deleted{+(-1)^{n-1}}\handwritten{+}\inserted{\handwritten{(-1)^{\unsure{n-1}}}}\operatorname{Trace}\operatorname{Frob}_{H^{n-1}(X_{\bar t}}\\
&\deleted{\unsure{+(-1)}}\handwritten{+(-1)^{\unsure{n-1}}}\operatorname{Trace}\operatorname{Frob}_{{\unsure{\mathbb{E}|\deleted{H}E|E}}^{n-1}(Y_s,\underline{Q}_{\handwritten{\ell}})}\,,
\end{aligned}
\]
\margin{\handwritten{$q=\operatorname{card}k(s)$}}
où tous les frobenius sont relatifs au corps de base $k(\deleted{\unsure{t}}s)$. Ici encore, quitte à remplacer $k_{\deleted{\unsure{t|s}}}$ par une extension finie, on peut supposer que les valeurs propres de\handwritten{s} frobenius (re\deleted{f}latif\handwritten{s} à $k$) opérant sur les $H^i(X_{\bar k},\underline{Q}_{\handwritten{\ell}})$ \inserted{\handwritten{$(0\leq i\leq n-2)$}} so\deleted{u}it égales à 0 ou 1 mod $p$ \inserted{\handwritten{(en tant qu'entiers algébriques)}}, de sorte que mod $p$ \deleted{les} la contribution des termes de la première ligne soit indépendante d\deleted{\unsure{u}}e choix de $\deleted{\unsure{t}}s$, \inserted{soit $c$.} Donc si les valeurs propres de frobenius opérant sur les $E^{n-1}(Y_s,\underline{Q}_{\handwritten{\ell}})$ \deleted{\unsure{étaient}}\inserted{sont} toutes divisibles par $p$, on en conclut\deleted{\unsure{ait|alors} que pour tout point fermé} que $\operatorname{card}Y_{s}(k(t))\equiv c\quad\bmod p$, et la même relation serait également valable, pour la même raison, si on remplace $k(t)$ par une extension finie $k'$ quelconque. \deleted{Or on peut d'après} Il suffit donc de voir qu'on peut trouver toujours un point \deleted{fermé} $s$ de $D\handwritten{\cap}U$ \deleted{tel que le}\inserted{dans une}\deleted{nombre}\inserted{extension finie $k'$ de $k$, tel que} $\operatorname{card}Y_{\deleted{\unsure{t}}s}(k\handwritten{'}\deleted{(s)})\deleted{\unsure{x}}$ soit congru mod $p$ à un élément donné de $\underline{F}_p$ (qu'o choisi\unsure{\deleted{r}t|\deleted{t}r}a alors $\handwritten{\ne}c$). Or c'est ce que montre essentiellement 7.2, à cela près qu'on n'a pas été assez précis dans l'énoncé de ce corollaire puisqu'on peut imposer à $\deleted{\unsure{t}}s$ de provenir de $D\handwritten{\cap}U\,(k')$, où $D$ est une droite donnée \inserted{\handwritten{$\in U'(k')$ \deleted{\unsure{passant}}}} rencontrant $U$. Ce raffinement est contenu dans la démonstration de 7.2, compte tenu du fait que \deleted{le}par \deleted{B}Bertini, le groupe fondame tal\deleted{e} de $\handwritten{(}D\handwritten{\cap}U\handwritten{)_{k'}}$ s'envoie \uline{sur} celui de $U$, donc s'envoie \uline{sur} $G_o$.
% archive sheet number: 20
